\section{The sampler representation}

Let $I$ be a finite or countable index set.  We isolate the consequence of the
OpenAI sampler that is used throughout the paper.

\begin{proposition}[OpenAI path representation]\label{prop:oai-input}
Suppose that a probability law $\nu$ on $\{0,1\}^I$ satisfies the hypotheses of
the Brownian sampling proposition in \cite[Section~4]{OpenAI2026FUSF}.  Let
$W=(W_i)_{i\in I}$ be independent standard Brownian paths, let $Z(W)$ be the
Borel solution of the sampler's deterministic inverse equation
\begin{equation}\label{eq:inverse-drift}
 Z_i(W)(t)=W_i(t)+\int_0^t b_i(s,Z(W)(s))\,\dd s,
\end{equation}
and let $\Psi(W)$ be its binary slope readout.  There exist a random configuration
$X\sim\nu$ and independent standard Brownian motions $(B_i)_{i\in I}$,
independent of $X$, such that
\begin{equation}\label{eq:oai-path-law}
  \bigl(Z(W),\Psi(W)\bigr)
  \ \stackrel{\mathrm d}{=}\
  \bigl((tX_i+B_i(t))_{i\in I,\,t\ge0},X\bigr).
\end{equation}
\end{proposition}

\begin{proof}
The cited proposition proves that independent Brownian input is transformed by
$Z$ into the complete joint path law of the reference observations
$Y_i(t)=tX_i+B_i(t)$.  Its slope readout is a Borel function of the path field,
and, in the reference model, that readout equals $X$ almost surely.  Applying
the same Borel readout to two path fields with the same law gives
\eqref{eq:oai-path-law}.  See the proof of the proposition entitled
``Sampling from independent Brownian paths'' in
\cite[Section~4]{OpenAI2026FUSF}.
\end{proof}

For the free uniform spanning forest, the OpenAI preprint verifies the required
uniform response estimate with Lipschitz constant $1/2$
\cite[Section~4]{OpenAI2026FUSF}.  Nothing in our argument below otherwise uses
spanning-forest structure.

In view of Proposition~\ref{prop:oai-input}, it is enough to work in the
reference model
\begin{equation}\label{eq:reference-model}
  Y_i(t)=tX_i+B_i(t),
\end{equation}
where $X$ has an arbitrary law on $\{0,1\}^I$ and the $B_i$ are independent of
$X$.

\begin{definition}[Midpoint readout and last error]\label{def:last-error}
For $t\ge0$, define
\begin{equation}\label{eq:midpoint-readout}
  \widehat X_i(t)=\1\{Y_i(t)>t/2\}.
\end{equation}
The last-error time at coordinate $i$ is
\begin{equation}\label{eq:last-error}
  T_i=\sup\{s\ge0:\widehat X_i(s)\ne X_i\}.
\end{equation}
For a finite $A\subseteq I$, put
\begin{equation}\label{eq:finite-stabilization}
  T_A=\max_{i\in A}T_i.
\end{equation}
\end{definition}

The supremum convention in \eqref{eq:last-error} is deliberate.  At a final
threshold contact, the strict readout in \eqref{eq:midpoint-readout} can be
correct for target bit $0$ and incorrect for target bit $1$.  In either case,
all times strictly greater than $T_i$ are decoded correctly.
